This appendix documents the revisions made to Deliverable 1 in response to the feedback received on the first submission. Entries are grouped under the headings used in that feedback so that each comment can be traced to the corresponding change. The revised document is written under the extended allowance of 2\,000 words granted with the feedback.

\begin{table}[htbp]
\centering
\small
\begin{tabular}{p{4.0cm} p{7.2cm} p{3.4cm}}
\toprule
\textbf{Feedback point} & \textbf{Change made} & \textbf{Where} \\
\midrule
Domain Description before Purpose and Scope
  & Section order inverted.
  & Domain Description \\
Intended Audience shortened or integrated
  & Standalone section removed; content condensed into a single paragraph.
  & Domain Description, final paragraph \\
Significance and Impact after the project outline
  & Moved from third to last section of the body.
  & Significance and Impact \\
Dedicated Data section
  & Added; test signals and datasets no longer described across several sections.
  & Data and Evaluation \\
State the intended mathematical formulation
  & New section: tokenisation operator, phase-locking conditions, locked
    frequency class, operational definition of structural aliasing.
  & Mathematical Formulation, Eqs.~(1)--(8) \\
Reformulate the hypotheses explicitly
  & Objectives replaced by H1--H3, each with a predicted direction, the metric
    used and an explicit refutation criterion.
  & Aliasing implications \\
Make the probing experiments precise
  & New section stating the probed representations, the metric, the probe
    points compared and the controls.
  & Probing Methodology \\
Simplify the experimental pipeline
  & Staged as: default checkpoint on clean sinusoids, then the retrained
    $(P,S)$ grid, then complex synthetic signals, then the photovoltaic data.
  & Data and Evaluation; Methodology \\
Clarify whether stride changes require retraining
  & Retraining shown to be necessary for both $P$ and $S$; the training recipe,
    its departures from the published one and their reasons are tabulated.
  & Methodology, Tables~1--2 \\
Explore strides on both sides of the default
  & Admissibility constraint stated and the symmetric design built within it.
  & Methodology \\
Describe the planned priors
  & Likelihoods, link function and all priors stated, with justification.
  & Bayesian Analysis, Eqs.~(9)--(10) \\
\bottomrule
\end{tabular}
\caption{Summary of the feedback received on the first submission and of the
corresponding revisions.}
\label{tab:changelog}
\end{table}

\subsection{Structure and organisation}
\begin{itemize}
  \item The order of the opening sections was inverted so that the application
        context precedes the statement of purpose. The document now opens with
        Domain Description, followed by Purpose and Scope.
  \item The standalone \emph{Intended Audience} section, previously two
        paragraphs, was removed. Its content was condensed into a single
        paragraph closing the Domain Description.
  \item \emph{Significance and Impact} was moved to the end of the body, after
        the reader has seen the formalisation, the data, the experimental
        design and the risk assessment.
  \item A dedicated \emph{Data and Evaluation} section was added. It now holds
        the synthetic sine sweep, the sinusoidal-mixer set, the KernelSynth
        contingency and the photovoltaic dataset, which in the first submission
        were split between \emph{Purpose and Scope} and \emph{Project Outline}.
  \item The single \emph{Project Outline} section was replaced by separate
        sections for the formalisation, the data, the training procedure, the
        probing protocol, the statistical analysis and the risk assessment.
  \item An appendix was added for the domain formalisation, so that the
        ontology does not consume the word allowance of the body.
\end{itemize}

\subsection{Mathematical formulation and hypotheses}
\begin{itemize}
  \item A \emph{Mathematical Formulation} section was added. The first
        submission announced the formalisation as a planned activity but did
        not outline it.
  \item The patch-based tokenisation operator $\mathcal{T}_{P,S}$ is defined
        explicitly as a map from the sampled signal to a sequence of
        vector-valued tokens, Eq.~(1).
  \item The signal model and the Nyquist constraint under which the phenomenon
        is studied are stated, Eq.~(2).
  \item The stride-translation derivation is given in full: the phase
        accumulated over one stride, Eq.~(3), the alignment condition on that
        accumulation, Eq.~(4), and the resulting exact stride-period
        invariance, Eq.~(5). The analogous condition for a shift of one patch
        length follows as Eq.~(6).
  \item The two conditions are collected into the phase-locking frequency class
        $\mathcal{F}_{\mathrm{lock}}$, Eq.~(7), expressed as a function of
        $f_s$, $P$ and $S$.
  \item It is now stated explicitly that self-repetition of consecutive patches
        is a property of a single signal and is not by itself evidence that two
        distinct signals become indistinguishable. The derived invariance and
        the empirical claim are therefore kept separate, and structural
        aliasing is given an operational, falsifiable definition at the level
        of the learned representations, Eq.~(8).
  \item The three items previously listed as objectives were reformulated as
        hypotheses H1, H2 and H3. Each now states a predicted direction, the
        quantity in which the effect should appear, the comparison against
        which it is measured, and the observation that would refute it.
  \item The ontology, previously a single bullet, is set out in
        Appendix~\ref{sec:appendixA}. The generation-state axioms are given in
        full; the anomaly, aliasing-event and agent-action axioms are described
        but not yet formalised, and are submitted at this stage for feedback on
        the modelling approach before the remaining rules are written.
\end{itemize}

\subsection{Experimental methodology}
\begin{itemize}
  \item A \emph{Probing Methodology} section was added, replacing the single
        sentence of the first submission.
  \item The probed representation is now specified: the 256-dimensional
        $[\mathrm{REG}]$ vector captured after each encoder block and each
        decoder block, together with the pre-projection input and the projected
        quantile output. The reason for probing this per-block vector rather
        than the flattened patch sequence is given.
  \item The first submission probed only after each decoder block. Probing the
        encoder as well allows the stage at which frequency information is lost
        to be located rather than merely detected.
  \item The metric is specified as a prequential MDL codelength, reported as a
        normalised codelength saving so that probe points of different
        dimensionality are compared on the same footing, together with the
        classification tasks on which it is computed and two controls, one with
        shuffled labels and one with an untrained network of identical
        architecture.
  \item A dense per-frequency accuracy profile was added, with deterministic
        non-redundant phase sampling in place of independently drawn random
        phases.
  \item The pipeline was simplified and ordered: the theoretical predictions
        are first tested on clean synthetic sinusoids with the released
        checkpoint, the architectural blind spots are then localised across the
        retrained configurations, and the photovoltaic dataset is used last, as
        an application case.
\end{itemize}

\subsection{Patch stride experiments}
\begin{itemize}
  \item It is now stated that any change to $P$ or $S$ requires retraining from
        scratch, with the reason for each: a change of $P$ resizes the
        patch-embedding block, and a change of $S$ alters the inter-patch
        dynamics the model has learned. The first submission announced a sweep
        over $S \in [8,16]$ without addressing this point.
  \item Table~1 reports the configurations retrained for this submission,
        with overlap, number of patches over the context window, step budget
        and wall-clock time.
  \item Table~2 compares every training hyperparameter against the published
        Chronos recipe and isolates the only departures from it, each with the
        constraint that motivates it.
  \item It is stated that the complete Chronos-Bolt training recipe is not
        publicly released, so the training loop is a budget-matched adaptation
        of the published Chronos-T5 recipe rather than an exact reproduction.
        The reduced step budget is applied uniformly to every configuration,
        including the default one, so that comparisons across stride values are
        not confounded by training budget. The released checkpoint is used only
        as an external forecast-quality reference.
  \item The exclusion of $S=4$ is justified rather than left implicit.
  \item The context window of Chronos-Bolt-Tiny is corrected to $T=2048$
        samples; the first submission reported $T=512$ tokens.
  \item Preliminary training runs are reported, with the hardware used and the
        total wall-clock time. Figure~1 is included as evidence that the
        proposed configurations train to a usable state within the available
        budget; it is not offered as an analysis of the phenomenon, which
        belongs to the later deliverables.
\end{itemize}

\subsection{Bayesian analysis}
\begin{itemize}
  \item A \emph{Bayesian Analysis} section was added, replacing the single
        sentence of the first submission.
  \item The analysis is now split into a relative and an absolute part. The
        relative part pairs each locked frequency with matched controls either
        side of it and reduces the pair to a logarithmic contrast, Eq.~(9), so
        that the localised effect is separated from the overall trend in
        forecast quality with frequency.
  \item The likelihood for the contrasts is stated, and a heavy-tailed choice
        is adopted to limit the influence of individual configurations. The
        prior on the population effect is given together with the reason for
        centring it at zero.
  \item The absolute part is specified as a Gamma regression with a log link,
        Eq.~(10), with the priors on the lock coefficient and on the dispersion
        parameter stated explicitly.
  \item The planned inference is stated: credible intervals and the posterior
        probability of an effect exceeding a pre-registered size, rather than
        point estimates alone.
\end{itemize}